\documentclass[10pt,a4paper]{article}
\usepackage[margin=1.6cm]{geometry}
\usepackage[T1]{fontenc}
\usepackage[utf8]{inputenc}
\usepackage{lmodern}
\usepackage{enumitem}
\usepackage{booktabs,tabularx,array}
\usepackage[hidelinks]{hyperref}
\usepackage[backend=biber,style=numeric,sorting=none,maxbibnames=3,doi=false,url=false,isbn=false,eprint=false,giveninits=true]{biblatex}
\addbibresource{refs.bib}
\setlist{nosep,leftmargin=1.2em,itemsep=1pt}
\setlength{\parindent}{0pt}
\setlength{\parskip}{3pt}
\usepackage[compact]{titlesec}
\titlespacing*{\section}{0pt}{6pt}{2pt}
\titlespacing*{\subsection}{0pt}{4pt}{1pt}
\newcommand{\adopt}{\textsc{[adopt]}}
\newcommand{\contrast}{\textsc{[contrast]}}
\newcommand{\cit}{\textsc{[cite]}}
\newcommand{\open}{\textsc{[open]}}
\usepackage{xcolor}
\usepackage{soul}
\sethlcolor{yellow!40}
\newcommand{\q}[1]{\hl{\footnotesize\itshape #1}}
\renewcommand{\bibfont}{\scriptsize}
\begin{document}

\begin{center}
  {\large\bfseries Literature summary for the supervisor meeting}\\[2pt]
  {\small Visual papers first (InfoSeek, Encyclopedic-VQA, M$^3$-VQA,
      ReasonVQA, OVEN, KVQA, RETINA, the KB-VQA audit, LLM-RA, SAGE-VQA,
    PMSR, KnowCoL), then the text-only papers
    (SPINACH, GRASP, QAWiki, GRISP, the WikiKGQA report, KGVoyager).
    Tags: \adopt\ idea we take, \contrast\ how we differ,
    \cit\ background, \open\ undecided.
    Numbers are from the papers (section/table in brackets).
  \q{Highlighted} = the open thesis question the point bears on.}
\end{center}

% \section{Visual benchmarks at a glance}
% \small
% \begin{tabularx}{\linewidth}{@{}p{1.55cm}p{1.9cm}p{1.6cm}p{1.8cm}p{1.6cm}X@{}}
%   \toprule
%   & KB & Gold SPARQL & Entity named? & Multi-hop & Index /
% distractors; agent; oracle \\
%   \midrule
%   InfoSeek \cite{chen2023can} & Wikidata (2022-10) + Wikipedia
% check & no & no (hypernym) & no (1 triple) & 100K Wikipedia KB
% incl.\ frequent entities; no agent; oracle entity \\
%   M$^3$-VQA \cite{ma2026m3} & Wikidata links, Wikipedia evidence &
% no & no (category) & yes, 3.1 avg & KB + $\sim$5 text-similar hard
% negatives per entity; planner/executor agent; entity-name oracle \\
%   ReasonVQA \cite{tran2025reasonvqa} & Wikidata & no (QID route
% only) & no (class) & yes, 1--3 & none (random MC answers); no agent; none \\
%   OVEN \cite{hu2023open} & Wikipedia (6M) & no & no (intent) & no &
% whole KB is the candidate set; no agent; n/a \\
%   KVQA \cite{shah2019kvqa} & Wikidata (2018) & withheld & mixed &
% yes, $\le$3 & 18K/69K persons; no agent; faces-solved oracle \\
%   \textbf{Ours} & Wikidata & \textbf{yes} & no & yes, labelled &
% shared CLIP index, same-class distractors; ReAct agent; entity-named oracle \\
%   \bottomrule
% \end{tabularx}
% \normalsize

\section{Themes across the visual papers}
\begin{itemize}
  \item \textbf{Recognising the entity is the bottleneck.} Giving the
    model the correct entity (oracle) helps a lot:
    InfoSeek $\sim$20\% $\to$ $\sim$50\% \cite[\S\,5.2, Table 5]{chen2023can};
    M$^3$-VQA entity-name oracle 24.5 $\to$ 34.7 \cite[\S\,5.3, Table
    4]{ma2026m3};
    KVQA closed world 49.5 $\to$ 59.2 \cite[Table 4]{shah2019kvqa};
    Encyclopedic-VQA: giving the subject category lifts PaLM 19.7
    $\to$ 31.0, and Lens recovers the right article only 47.4\% of
    the time, with 82.3 vs.\ 20.7 accuracy for a correct vs.\ wrong
    retrieved section \cite[\S\,5.2, \S\,5.4, Tables
    5--6]{mensink2023encyclopedic}.
    On OVEN the best model reaches 20.2 harmonic mean, humans with a
    search engine 77.7 \cite[\S\,5.2, Table 2]{hu2023open}.
    \adopt\ supports our oracle / vision-only / index / both
    configurations \q{Which agent configurations do we evaluate?}
  \item \textbf{Accuracy falls as the candidate set grows.} OVEN
    CLIP2CLIP drops roughly linearly in log(candidates) from 20K to
    6M \cite[\S\,6, Fig.~5]{hu2023open};
    KVQA loses more than 10 points from 18K to 69K identities
    \cite[\S\,Approach, Table 4]{shah2019kvqa}.
    \adopt\ motivates varying the number of distractors \q{How many
    distractors?}
  \item \textbf{Nobody picks visually similar distractors.} OVEN,
    KVQA and InfoSeek use all other entities (InfoSeek adds frequent
    Wikipedia entities) \cite[\S\,4]{chen2023can};
    M$^3$-VQA samples 5 of the top-20 text-summary neighbours
    \cite[App.~D.3]{ma2026m3}; ReasonVQA only has random
    multiple-choice answers \cite[\S\,3.2.2]{tran2025reasonvqa}.
    \contrast\ our lean: entities of the same class (\texttt{instance
    of}) \q{How are distractors chosen?}
  \item \textbf{SPARQL builds the data but is never released.} Human
    templates per relation plus SPARQL for answers in InfoSeek (300
    relations), M$^3$-VQA, ReasonVQA (182 + 100 templates) and KVQA,
    whose queries are withheld on purpose \cite[S3]{shah2019kvqa}.
    ReasonVQA's Hugging Face data comes closest: main-entity QID,
    route of QIDs, last property.
    \contrast\ \q{Does an image + question + gold-SPARQL benchmark
    exist?} -- gold SPARQL is our difference.
  \item \textbf{The image can show a related entity, not the answer
    entity.} RETINA removes this ``visual shortcut'': the query image
    shows a related entity of the answer document.
    Single-image retrievers fall from 51.7 (InfoSeek) / 64.9 (E-VQA)
    to 20.8 / 8.6 Recall@5 on RETINA seen / unseen; adding images of
    related entities to the document gets MIMIR to 43.2 / 36.9
    \cite[Table 1]{lee2025breaking}.
    \adopt\ (leaning) as an option, with the related entity a Wikidata
    neighbour of the answer entity \q{Can the image show another
    entity?} \q{What role does the image play?} \q{How is difficulty defined?}
  \item \textbf{Popularity changes which approach wins.} End-to-end
    models win on tail entities, CLIP pipelines on popular ones
    \cite[\S\,5.2, Fig.~7]{chen2023can}; same pattern with 5K/2.5K
    page-view thresholds in OVEN \cite[Fig.~6]{hu2023open}.
    \adopt\ candidate difficulty axis and label \q{How is difficulty
    defined?} \q{Which labels does each question carry?}
  \item \textbf{Image-side tools: crop the entity before searching.}
    LLM-RA: cropping the key-entity regions raises InfoSeek
    Recall@5 from 39.8 to 47.3, and without cropping recall falls as
    the number of objects grows (41.6 $\to$ 37.1, 1--3 vs.\ 9+
    objects) \cite[Tables 3--4]{jian2024large}; M$^3$-VQA's agent
    already detects and crops \cite[\S\,5.4]{ma2026m3}; SAGE-VQA
    adds a VLM check of each crop before it is used for retrieval
    \cite[\S\,3.2, App.~A.7]{shen2026sage}. \open\ crop tool, maybe
    with one merged check, as an agent tool (idea liked, not decided)
    \q{Which tools does the
    agent get?} \open\ whole image vs.\ unchecked crops vs.\ checked
    crops as an ablation, only if we build the tool \q{Which ablations?}
    \q{Which agent configurations do we evaluate?}
  \item \textbf{Benchmarks contain flawed and underspecified items.}
    In InfoSeek 22\% of annotations are not supported by the evidence
    (E-VQA 1\%); 47\% (InfoSeek) and 59\% (E-VQA) of questions are
    ambiguous, mostly a missing attribute, time or place
    \cite[\S\,3.1--3.2]{ma2026identifying}. \adopt\ three review
    checks per item: evidence supports the answer, constraints give a
    unique answer, no answer leakage \q{How are questions collected and
    curated?} \q{How are gold answers defined and verified?}
    Adding one distractor entity to the query image drops InfoSeek
    R@1 from 43.5 to 14.7 / 20.4 (same / other category, Table 2).
    \open\ as a separate difficulty setting, not as index distractors
    \q{How is difficulty defined?} \q{How are distractors chosen?}
  \item \textbf{Distractors in the context hurt.} PMSR, E-VQA subset
    with the oracle text: 86.6\% accuracy, 78.7\% after adding 10
    retrieved image-text pairs, 64.3\% with 20 pairs
    \cite[App.~M, Table A8]{choi2025progressive}. \cit\ evidence that
    visually similar irrelevant images matter \q{How do related works
    choose distractors?}
  \item \textbf{No-image / no-knowledge baselines.} RETINA: a
    retriever trained without the question text does about as well on
    InfoSeek and E-VQA \cite[\S\,1, Fig.~2]{lee2025breaking}.
    M$^3$-VQA question-only 15.6 vs.\ 24.5 with the image \cite[Table
    4]{ma2026m3}; InfoSeek question-only PaLM 4.3, and 70.8\% of
    sampled OK-VQA questions need no search \cite[\S\,2, Table 4]{chen2023can};
    ReasonVQA is image + question without tools, GPT-4o 62.8
    open-ended \cite[Table 2]{tran2025reasonvqa}. \adopt\ our
    ``vision only'' baseline \q{Which baselines?}
  \item \textbf{They score answer text, we compare query results.}
    Aliases, $\pm$1 year, 10\% for numbers
    \cite[App.~A.3]{chen2023can}; set IoU \cite[\S\,4.3]{ma2026m3};
    exact/substring/semantic similarity, exact match low because of
    free text \cite[\S\,4.3]{tran2025reasonvqa}. \contrast\ \q{Which
    metrics do we report?} \q{Which KGQA metrics exist?}
  \item \textbf{Diversity of the question set is measured.}
    Encyclopedic-VQA reports variety (unique bigrams, 257.9K) and
    disparity (mean pairwise Sentence-BERT cosine distance, 0.833)
    \cite[\S\,4.5, Table 3]{mensink2023encyclopedic}. Variety grows
    with dataset size; disparity could be computed on our small
    curated set. \open\ \q{How large should the dataset be?}
  \item \textbf{Taxonomies and difficulty axes.} KVQA: 10
    non-exclusive categories, oracle accuracy from Boolean 75.1 to
    subtraction 40.5 \cite[Fig.~2, Table 5]{shah2019kvqa}; M$^3$-VQA:
    parallel vs.\ sequential, hops and entities \cite[\S\,3]{ma2026m3};
    InfoSeek: answer types Time/Numerical/String, Time hardest
    \cite[Table 6]{chen2023can}; Encyclopedic-VQA: templated,
    automatic, multi-answer and two-hop via a bridge entity, two-hop
    clearly harder (PaLM 22.8 vs.\ 48.8, \S\,4.3, \S\,5.6)
    \cite{mensink2023encyclopedic}; ReasonVQA: by property, 3-hop
    clearly harder, 1-hop barely easier than 2-hop
    \cite[\S\,4.4]{tran2025reasonvqa}. \adopt\ tags per question
    \q{Which labels does each question carry?} \q{Which question
    taxonomy to use?}
\end{itemize}

\section{Per paper}
\subsection*{InfoSeek \cite{chen2023can}}
Entity-free questions about a depicted Wikipedia entity: 8.9K human +
1.35M Wikidata-template questions (Table 2).
\begin{itemize}
  \item Templates with a hypernym placeholder (``this car''); human
    split rejects questions that copy the entity name (\S\,3.2,
    App.~A.1). \adopt\ \q{How are questions collected and curated?}
  \item Answer kept only if found in the entity's Wikipedia article
    (97\%), and subsampling by (relation, category) to avoid skew
    (App.~A.2). \adopt\ \q{How are questions collected and curated?}
  \item With-KB protocol: image + question embedding vs.\ (Wikipedia
    image, title) in a 100K KB; oracle entity gives 52.0 vs.\ 20.1
    (Table 5). \contrast\ our index is image-only with a QID payload
    \q{Which embedding model?} \q{What is stored per entity?}
  \item Single triple per question, no SPARQL, no agent,
    Google-internal models (\S\,3.2, \S\,5). \contrast\ \q{Does an
    image + question + gold-SPARQL benchmark exist?}
\end{itemize}

\subsection*{Encyclopedic-VQA \cite{mensink2023encyclopedic}}
1M image-question-answer triplets (221K question pairs, up to 5
images each) about species (iNat21) and landmarks (GLDv2); the
knowledge base is 2M Wikipedia articles, answers are tied to a
supporting section (\S\,4, Tables 1--2).
\begin{itemize}
  \item The subject is named only by its super category (``this
    monastery''); the same question occurs with different answers, so
    the image is needed (\S\,4). \adopt\ supports our entity-free questions.
  \item Two-hop questions chain two single-hop questions through a
    bridge entity with its own page; PaLM chains them and a second
    prompt validates them (\S\,4.3). \open\ chaining as an idea for
    Wikidata paths, not decided \q{What is the maximum hop depth?}
    \q{Which question taxonomy to use?}
  \item Staged oracle settings (subject category, KB article, KB
    section) versus real retrieval with Google Lens (\S\,5.1--5.3).
    \adopt\ only the subject oracle \q{Which agent configurations do
    we evaluate?}
  \item Accuracy uses BEM; exact match agrees with humans in only
    68\% of cases (App.~B). \contrast\ \q{Which metrics do we report?}
  \item Wikipedia text KB, no SPARQL or graph, no released image
    index, no agent (\S\,4--5). \contrast\ \q{Does an image +
    question + gold-SPARQL benchmark exist?}
\end{itemize}

\subsection*{M$^3$-VQA \cite{ma2026m3}}
13,125 questions, avg.\ 3.1 hops and 2.8 entities; images may show
several entities (Table 2).
\begin{itemize}
  \item Closest to our design: controlled KB with hard negatives,
    oracle levels (entity name / sentence / section),
    planner--executor--solver agent with detection, text and image
    retrieval (\S\,5.4). \adopt\ \q{Which tools does the agent get?}
    \q{Which agent configurations do we evaluate?}
  \item Agentic retrieval beats one-shot retrieval for all 7 solvers,
    e.g.\ 22.5 / 25.7 / 31.2 (original / heuristic / agentic,
    Qwen2.5-VL-7B, Table 5); caveat: GPT-4o is the planner.
  \item Accuracy falls with hops/entities under the sentence oracle,
    smaller fall in the original setting (image + question) (\S\,5.5,
    Fig.~2). \adopt\ \q{What is the maximum hop depth?} \q{How is
    difficulty defined?}
  \item Entity embedding = CLIP article summary + image embedding,
    better than either alone (App.~F.4). \cit\ \q{Which embedding model?}
  \item Released data has no QIDs and no SPARQL; \contrast\ \q{Does
    an image + question + gold-SPARQL benchmark exist?}
\end{itemize}

\subsection*{ReasonVQA \cite{tran2025reasonvqa}}
4.2M questions on 598K images (Visual Genome, Google Landmarks v2);
1--3 hops from nested templates (\S\,3.2.1).
\begin{itemize}
  \item Visual Genome scene-graph relations as sub-clauses add visual
    context and lower scores (\S\,3.2.1, \S\,4.4). \open\ \q{Which
    visual-KG papers inspire us?} \q{Is the image always the single
    subject?} \q{How is difficulty defined?}
  \item Hops capped at 3 for grammar; only 5.4K of 4.2M questions are
    3-hop (App.~1.2, Table 2). \cit\ \q{What is the maximum hop depth?}
  \item Released rows: QID, route, last property only; sample of 6.7K
    rows all landmarks; AGPL-3.0; repository empty at check time.
    \open\ \q{How are questions collected and curated?} -- rebuilding
    gold SPARQL for about 20 items is a possible later check, not
    planned; licences of Commons images unchecked.
\end{itemize}

\subsection*{OVEN \cite{hu2023open}}
Image + intent query $\to$ one of 6.06M Wikipedia entities;
single-step recognition (\S\,2--\S\,3).
\begin{itemize}
  \item Annotators rewrote queries so no other object fits and the
    entity is not leaked; 8.4\% of entity links wrong, 10.3\% of
    answers not visible (App.~10.2). \adopt\ \q{How are questions
    collected and curated?} review step.
  \item Vanilla CLIP ViT-L/14 is weak at 6M entities (Table 2),
    CLIP-style retrieval relatively better on tail entities (Fig.~6).
    \cit\ \q{Which embedding model?}
  \item Error analysis of 100 examples: wrong-but-relevant 23\%, too
    generic 15\%, misunderstood query 7\% (Table 3). \open\ \q{Which
    ablations?} -- small error analysis in the thesis?
  \item Category skew (animal 19.4\%, plant 16.0\%, building 13.8\%,
    Fig.~3). \cit\ \q{Which part of Wikidata is well connected?} only as
    a reminder that the entity pool is skewed by category.
\end{itemize}

\subsection*{RETINA / MIMIR \cite{lee2025breaking}}
Wikipedia KB-VQA where the image shows a related entity; 120K
generated training questions, 2K human-curated test (1K seen, 1K
unseen) (\S\,3.2).
\begin{itemize}
  \item Leakage controls: demonstrative instead of the entity name,
    BM25 top-5 check on the question (37\% $\to$ 25\% retrievable,
    rest dropped), separate query and document image pools with a
    perceptual-hash check, human removal of questions solvable
    without the image (App.~A). \adopt\ conditions 1 and 4 as review
    criteria, 2 and 3 automatable \q{How are query images selected
    and leakage prevented?}
  \item Hard distractors: documents whose main entity matches the
    query image; MuKA retrieves them more often than the gold
    document on unseen (\S\,5.2, Fig.~5). \cit\ \q{How are distractors chosen?}
  \item Pipeline: LLM one-hop graph, query + qualifying entity for a
    unique answer, only 29.2\% of 6K sampled questions survive
    curation, mostly because of several valid answers (App.~A.3).
    \open\ the 2K test set is a likelier source of examples than the
    pipeline \q{How are questions collected and curated?}
  \item Wikipedia text KB, no SPARQL, no agent (retrievers +
    VILA-13B). \contrast\ \q{Does an image + question + gold-SPARQL
    benchmark exist?}
\end{itemize}

\subsection*{KVQA \cite{shah2019kvqa}}
183K questions on 24.6K images of named persons, Wikidata dump
2018-05-05, up to 3 hops (Table 2, \S\,Knowledge Graph).
\begin{itemize}
  \item Ten non-exclusive question categories as tags (Fig.~2).
    \adopt\ \q{Which labels does each question carry?}
  \item Paraphrased test questions cost about 17 points (49.5 $\to$
    32.0, Table 4). \open\ \q{How are questions collected and
    curated?} paraphrasing?
  \item Persons by face only, 2019 baselines, no agent, queries
    withheld. \contrast\ \q{Does an image + question + gold-SPARQL
    benchmark exist?}
\end{itemize}

\subsection*{KB-VQA audit \cite{ma2026identifying}}
Audits and repairs InfoSeek and E-VQA, then composes distractor
entities into the query image (\S\,3--6).
\begin{itemize}
  \item Unsupported annotations 22\% / 1\%, ambiguous questions
    47\% / 59\% (InfoSeek / E-VQA); fixing changes method rankings
    (\S\,3, \S\,6.2). \adopt\ the three checks as review criteria
    \q{How are questions collected and curated?} adapted to the frozen
    snapshot and the gold SPARQL result \q{How are gold answers
    defined and verified?}
  \item Composite query images with one extra entity: intra- and
    inter-category variants, InfoSeek R@1 43.5 $\to$ 14.7 / 20.4
    (Table 2). \open\ possible separate difficulty setting; RETINA or
    M$^3$-VQA multi-entity images fit our small set better than
    automatic composites \q{How is difficulty defined?} \q{What role
    does the image play?} \contrast\ not an index-distractor result
    \q{How are distractors chosen?}
  \item Wikipedia evidence, no executable query. \contrast\
    \q{Does an image + question + gold-SPARQL benchmark exist?}
\end{itemize}

\subsection*{LLM-RA \cite{jian2024large}, SAGE-VQA \cite{shen2026sage}}
Both put an image-side step before retrieval, not datasets.
\begin{itemize}
  \item LLM-RA: an LLM picks the key entities from caption and
    question, a grounder crops them, each crop is encoded separately.
    Without the crops Recall@5 falls 47.3 $\to$ 39.8 (InfoSeek), and
    recall drops as the object count grows (Tables 3--4). \cit\
    cropping helps; we take the idea via M$^3$-VQA's detect step, not
    its trained retriever \q{Which tools does the agent get?}
  \item SAGE-VQA: only VLM-verified crops go to image retrieval; failed
    crops stay visible with a flag (\S\,3.2, App.~A.3). \open\ one
    merged check that sees crop and whole image
    \q{Which tools does the agent get?}
  \item Caveat: SAGE-VQA has no ablation of the verification itself
    (``+EDAR'' also adds a readiness loop, Table 5). \open\ our own
    ablation: whole image vs.\ unchecked vs.\ checked crops
    \q{Which ablations?} Its memory bank is out of scope.
\end{itemize}

\subsection*{PMSR \cite{choi2025progressive}}
Training-free retrieve--condense--requery loop (\S\,3); E-VQA BEM
37.1 $\to$ 47.1 over iterations (Table 5).
\begin{itemize}
  \item \open\ an optional iterative mode for our agent, not committed
    \q{Which tools does the agent get?} It is a fixed workflow, not
    ReAct \q{Build an own agent or build on GRASP?}
  \item With oracle text, 10 retrieved image-text pairs cut accuracy
    86.6 $\to$ 78.7\% (App.~M, Table A8), see distractor theme.
\end{itemize}

\subsection*{KnowCoL \cite{zhou2025seeing}}
Trained OVEN dual encoder aligned with a TransE embedding of a Wikidata
subgraph; 38.8 harmonic mean vs.\ 11.4 for CLIP2CLIP (Tables 1--2).
\contrast\ the graph is used only in training and needs 4--8 large
GPUs, we train nothing \q{Which embedding model?} \q{What is stored
per entity?}

\section{Text-only agent papers}
\begin{itemize}
  \item \textbf{SPINACH} \cite{liu2024spinach}: tools
    \texttt{search\_wikidata}, \texttt{get\_wikidata\_entry},
    \texttt{get\_property\_examples}, \texttt{execute\_sparql};
    removing search drops F1 46.4 $\to$ 25.3 (Table 4); row-major F1;
    zero-shot high-level prompt. \adopt\ \q{Which tools does the
    agent get?} \q{Which metrics do we report?} \q{Which ablations?}
    (tool removal, error analysis)
  \item \textbf{GRASP} \cite{walter2025grasp}: execute-only tool set
    is worst, simpler tools beat flexible ones (Table 6); Qwen2.5-72B
    68.7 vs.\ GPT-4.1 75.3 on WWQ (Table 5). \cit\ \q{Own agent or
    build on GRASP?} (leaning: own agent, GRASP as a possible baseline)
    \q{Which baselines?}
  \item \textbf{QAWiki} \cite{loustaunau2025qawiki}: text-only,
    handcrafted Wikidata dataset, 526 questions (518 with SPARQL);
    answer-type categories (entity, entity-set, numeric, temporal,
    boolean, other; \S\,3.1) and quality tags (ambiguous, subjective,
    incomplete, no ties; \S\,3.5). SPINACH's row F1 is reused
    unchanged (\S\,4.1). GPT-4o F1 0.164 $\to$ 0.406 when given the
    linked entities (Table 2). \adopt\ a second oracle with the
    subject QID \q{Which agent configurations do we evaluate?};
    \adopt\ exclude tagged items \q{How are questions collected and curated?}
  \item \textbf{GRISP} \cite{walter2026grisp}: fine-tuned 7B model
    that fills SPARQL skeletons by retrieve-and-re-rank; needs training
    pairs, wins only where training data exists (Table 3). \cit\ not a
    baseline since we do not train \q{Which baselines?}
  \item \textbf{WikiKGQA report} \cite{walter2026grasp}: GRASP and
    GRISP in a Wikidata text-to-SPARQL challenge on a fixed dump served
    from a local QLever (\S\,2--3). GRASP with GPT-5.4 reaches 85.0
    macro QALD-F1 (English, no mentions); an off-the-shelf Claude Code
    agent matches it at about \$55 vs.\ \$9 (Table 2); GPT-5.4 86.8
    vs.\ mini 78.1 validation F1, so the backbone matters (Table 1).
    \open\ ideas only, not decided: agent-written exploration notes and a shape
    index of class properties as a schema tool \q{Which tools does the
    agent get?} \cit\ \q{Build an own agent or build on GRASP?}
    \q{Which CampusAI model and setup?} \q{Which Wikidata endpoint and
    snapshot?} \contrast\ the Claude Code comparison uses paid APIs,
    out of scope (D-010) \q{Which baselines?} QALD-F1 (harmonic mean
      of macro P and R, empty answer
    gets P=1) is not the SPINACH row F1 and is not reported
    \q{Which KGQA metrics exist?} \q{Which metrics do we report?}
    The authors name Wikidata's presence in LLM training data as a
    threat \q{What are the research questions and hypotheses?}
\end{itemize}
\textbf{KGVoyager} \cite{wisam2026kgvoyager}: ReAct agent for any RDF
endpoint that indexes classes instead of all entities.
\begin{itemize}
  \item F1 64.6 vs.\ GRASP 60.9, cost $-$38.3\%, time $-$23.8\% on
    four small benchmarks (88--108 questions) (Table 3). Not usable
    on Wikidata (no \texttt{rdf:type} convention; Limitations, Table 2).
  \item Tools worth borrowing: hybrid (dense + BM25) search, dropping
    low-score results against ``context rot'' ($\delta$=0.125),
    property-usage exploration (\S\,3.1--3.2). \adopt\ \q{Which tools
    does the agent get?} \q{Which CampusAI model and setup?} Hybrid
    beats keyword-only by 1.1--26.5 F1 (Table 4).
  \item Failure analysis of 156 complete failures: 36\% agent gave up
    (small models, long contexts), 19\% concept resolution, 18\%
    answer representation (App.~F). \cit\ \q{How are gold answers
    defined and verified?} \q{Which CampusAI model and setup?} Also
    reports cost, iterations, time (Table 3). \adopt\ \q{Which
    metrics do we report?}
\end{itemize}

\section{How we differ}
\q{Does an image + question + gold-SPARQL benchmark exist?}
None of the visual benchmarks and methods releases gold SPARQL, scores by
executing a query, evaluates an agent that queries the KG
(M$^3$-VQA's agent retrieves Wikipedia text/images only), or covers
non-person entities with entity-free questions
\cite{chen2023can,mensink2023encyclopedic,ma2026m3,tran2025reasonvqa,hu2023open,shah2019kvqa,lee2025breaking,ma2026identifying,choi2025progressive,zhou2025seeing}.
The text-only Wikidata benchmarks have gold SPARQL but no images
\cite{loustaunau2025qawiki,walter2026grasp}.

\section{Questions}
\begin{enumerate}
  \item Use ReasonVQA items (QID routes exist, landmarks only, AGPL)
    as a cheap source of image + question pairs, or only as a fallback
    to the LLM-suggested, manually curated pipeline (D-016)?
  \item Distractors: same-class siblings vs. M3-VQA text neighbours
    vs. an InfoSeek-style frequent-entity pool (leaning: same class,
    others as additions); sweep the distractor count?
  \item Entity popularity (page views) as difficulty axis and label;
    state 'vision-only does better on popular entities' as a hypothesis?
  \item QID oracle as second oracle next to the entity name (decided in
    AG-05, QAWiki-style); drop the SPARQL-triples oracle?
  \item Report only execution-based row F1 (plus cost, iterations); is a
    label-text metric needed for comparability with the visual papers?
    (QALD-F1 not planned)
  \item Tags (KVQA categories, hops, answer type) and hop cap of 3
    vs. M3-VQA's 4+?
  \item Images of related entities (RETINA-style, \q{Can the image
    show another entity?}) as an option for harder questions (leaning
    yes); the image-only check is already planned for the pilot.
  \item Diversity of the dataset: report disparity (Encyclopedic-VQA
    style) for our small curated set, or skip?
  \item Open model on CampusAI: cap tool output length given the
    'agent gave up' failures?
  \item Crop tool for the image index (LLM-RA, M3-VQA, SAGE-VQA): build
    it with one merged check, and ablate whole image vs.\ unchecked
    vs.\ checked crops, or skip (needs image input on CampusAI)?
  \item Confirm the audit's three checks (evidence support, unique
    answer, no leakage) as acceptance criteria (tentative); composite
    query images only as an optional separate difficulty setting, or
    prefer RETINA/M3 multi-entity images?
  \item Iterative PMSR-style mode: optional switch, not committed?
\end{enumerate}

\printbibliography
\end{document}
